% ============================================================
%  1. Introduction
% ============================================================
\section{Introduction}
\label{sec:intro}

Modern NER pipelines overwhelmingly rely on contextual language models such
as BERT~\citep{devlin2019bert}, RoBERTa~\citep{liu2019roberta}, or
DeBERTa~\citep{he2021deberta}.  These models preprocess text through fixed
subword vocabularies---WordPiece or Byte-Pair Encoding (BPE)---and their
performance is conditioned on the assumption that input tokens appear in a
form close to those seen during pretraining.

This assumption routinely fails in practice for several reasons:

\begin{itemize}
  \item \textbf{User-generated content} (social media, support tickets, search
        queries) contains frequent spelling errors.
  \item \textbf{OCR pipelines} introduce substitution, insertion, and deletion
        noise at the character level.
  \item \textbf{Multilingual and code-switched} text produces tokens that fall
        outside any fixed vocabulary.
  \item \textbf{Adversarial inputs} exploit tokenizer blindspots to evade entity
        detection systems.
\end{itemize}

As a concrete example, when BERT encounters the query \emph{``Micorsoft
anounced a partership with Opnai''}, its WordPiece tokenizer fragments the
misspelled tokens into meaningless subword pieces, degrading NER recall for
the very entities that matter most downstream.

Three families of approaches attempt to address this fragility:
(i)~\emph{character-level models}~\citep{lample2016neural,ma2016end} that
lack contextual pretraining and are computationally expensive at inference;
(ii)~\emph{robustness augmentation}~\citep{pruthi2019combating} that
fine-tunes pretrained models on noisy text but remains bottlenecked by the
subword tokenizer at its input; and
(iii)~\emph{hybrid tokenization}~\citep{xue2022byt5,clark2022canine} that
partially decouples character and word representations, adding architectural
complexity.

We take a different path.  Our key observation is:

\begin{quote}
\emph{Character-level embeddings can serve as a universal, vocabulary-free
word representation if distilled carefully from a high-quality subword-based
teacher.}
\end{quote}

We propose \LightRet{}, which combines three components:

\begin{enumerate}
  \item \textbf{\RetVec{}}~\citep{sander2023retvec}: a compact, pretrained
        character-level embedder that maps any Unicode word to a
        $256$-dimensional vector without a vocabulary lookup, and whose
        representation space is explicitly regularised so that typographically
        similar words are close.
  \item A \textbf{three-stage progressive distillation pipeline} that transfers
        BERT's contextual knowledge into a lightweight
        BiGRU--Transformer backbone via an intermediate \RetBERT{} student.
  \item A \textbf{noisy-student NER fine-tuning} stage with a stochastic
        character-noise simulator and a novel \emph{dynamic BIO label projection}
        algorithm that handles word-splitting noise without loss of label
        integrity.
\end{enumerate}

\paragraph{Contributions.}
\begin{enumerate}
  \item A principled three-stage teacher--student distillation chain
        $\textsc{BERT} \!\to\! \RetBERT{} \!\to\! \LightRet{}$, where the final
        student uses a fundamentally different input representation
        (character-level \RetVec{}) than the teacher (subword-based BERT).
  \item A noisy-student NER training procedure featuring a character-level
        noise simulator and a dynamic BIO label projection algorithm that
        correctly handles merge and split events caused by space-insertion
        noise.
  \item A compact (${\sim}4$M parameter), vocabulary-free NER model
        deployable on CPU endpoints, achieving competitive clean-text F1 and
        superior noise robustness compared to BERT-base and CharBERT baselines.
  \item Comprehensive ablation studies isolating the contribution of each
        distillation stage, the noise augmentation, and the distillation loss
        component.
\end{enumerate}
